%% !TEX root = ../manual_gaolv5.tex
% Chapter 1.
%% Chapter 1: Interval arithmetic and floating-point numbers. Sources:
%% gaol.tex of f5bf180 141-257, 840-938, 1109-1150, 1913-1996, 2840-2846, 3493-3546,
%% 4070-4099, 5532-5546, 5574-5620, 5781-5788, 5881-5931;
%% doc/math-libraries.md (whole); doc/accuracy.md 1-47.
\chapter{Interval arithmetic and floating-point numbers}\label{chap:concepts}
\cindex{interval arithmetic}\cindex{floating-point numbers}%

A computer holds a real number as a floating-point number, a \code{double},
and a computation on doubles rounds each of its results: the quotient $1/3$,
the constant $0.1$ and most values of the elementary functions have no exact
representation, and a program that computes with them accumulates errors it
does not see. Interval arithmetic replaces each real number by two doubles
that enclose it, and each operation by an operation on intervals whose
result encloses every value the exact operation can take. The errors are
then not avoided but bounded: the result of an interval computation is a set
that contains the exact result, whatever the roundings along the way. GAOL
is a C++ library to perform arithmetic with floating-point intervals, and
this chapter gathers the notions on which the rest of the manual relies.

\medskip
In this chapter, we define the floating-point numbers and their roundings,
the floating-point intervals and the operations on them, and the vocabulary
of IEEE 1788-2015, the standard for interval arithmetic that GAOL v5
follows. The first section describes the doubles and the roundings
downward, upward and to nearest. The second defines the intervals, the
empty set, the unbounded intervals and the canonical intervals. The third
states the containment property and the outward rounding that enforces it,
with the property on which GAOL computes both bounds in a single rounding
direction. The fourth presents the dependency problem, which makes an
enclosure wider than the range. The fifth gives an overview of IEEE
1788-2015 as GAOL v5 follows it, and the sixth the idea of the relational
functions, which the standard calls reverse functions. The seventh explains
why the bounds of an elementary function depend on the mathematical library
that computes it, from the libm of the systems to CORE-MATH; the eighth
defines the three accuracy modes of the standard, tightest, accurate and
valid. The chapter ends with a discussion.

\section{Floating-point numbers and rounding}\label{sec:rounding-definitions}
%%------------------------------------------
\cindex{F@\FSet}\cindex{R@\RSet}\cindex{double}%

The numbers GAOL computes with are the \code{double}s of C++: the
floating-point numbers of the double format of IEEE
754~\cite{IEEE754:2008,Goldberg:01}, called binary64 since the 2008 revision
of the standard. A double is written on 64 bits: a sign, a significand of
53 bits, 52 of them stored, and an exponent of 11 bits; its value is
$\pm m\cdot2^{e}$ with $m$ an integer below $2^{53}$. We write \FSet{} for
the set of the doubles, the two infinities included, and \RSet{} for the set
of the real numbers. \FSet{} is finite and symmetric: $y\in\FSet$ if and
only if $-y\in\FSet$.

\medskip
The doubles are not spread evenly on the real line. Between two consecutive
powers of two, $2^{e}$ and $2^{e+1}$ with $-1022\leq e\leq1023$, they are $2^{52}$ in number, spaced by
$2^{e-52}$: this spacing is the \dfn{unit in the last place} (ulp) of the
doubles between these powers. The \dfn{normal numbers} are those of
magnitude between $2^{-1022}$, the smallest normal double, and
$\mathrm{MAX}=(2-2^{-52})\cdot2^{1023}\approx1.797693\cdot10^{308}$, the
largest double, which C names \code{DBL\_MAX}. Below $2^{-1022}$ lie the
\dfn{subnormal numbers}, the multiples of $2^{-1074}$, which fill the gap
between the smallest normal double and zero with the same spacing. The
bounds GAOL computes may be subnormal, and the modes of the processors that
flush the subnormal numbers to zero break them (\pxref{sec:ftz}).
\cindex{subnormal numbers}\cindex{unit in the last place}%

\medskip
Three kinds of values are special. The two zeros, $-0$ and $+0$, are equal in the
comparisons and stand for the same number; IEEE 1788-2015 and GAOL take them
for the same bound, and the sign a zero bound gets differs between the
builds of GAOL (\pxref{sec:zero-sign}). The two infinities, $-\infty$ and
$+\infty$, which an overflow or a division by zero produces, are the bounds
of the unbounded intervals. NaN (\emph{not a number}) is the result of an
invalid operation, as $0/0$ or $\sqrt{-1}$; it compares unequal to
everything, itself included, and GAOL holds the empty set as an interval
whose two bounds are NaN (\exref{sec:intervals}).
\cindex{NaN}\cindex{zero!sign of a bound}\cindex{infinity}%

\medskip
A real number that is no double has to be rounded. IEEE 754 defines four
rounding directions for the binary formats of the processors: to nearest, toward $-\infty$, toward $+\infty$ and
toward zero. The processor holds the direction in effect in a control
register, and every operation of its floating-point unit rounds its result
in that direction; a program starts in the rounding to nearest. Interval
arithmetic uses the first three, which we define now.
\cindex{rounding!direction}%

\begin{Def}[Round down, round up]\label{def:round-down-up}
\cindex{rounding!downward|hyperpagebf}\cindex{rounding!upward|hyperpagebf}%
\cindex{downward rounding ($\roundDn{\:}$)}%
\cindex{upward rounding ($\roundUp{\:}$)}%
\index[idx]{aaa@$\roundDn{\:}$|see{downward rounding}}%
\index[idx]{aab@$\roundUp{\:}$|see{upward rounding}}%
Given \RSet{} the set of real numbers and \FSet{} the set of floating-point
numbers (\code{double}), we have, for all $x\in\RSet$:
\begin{align*}
  \roundDn{x} &= \max\{y\in\FSet\mid y\leq x\},\\
  \roundUp{x} &= \min\{y\in\FSet\mid y\geq x\}.
\end{align*}
$\roundDn{x}$ is the rounding of $x$ \dfn{downward}, toward $-\infty$, and
$\roundUp{x}$ its rounding \dfn{upward}, toward $+\infty$.
\end{Def}

\begin{Def}[Round to nearest]\label{def:round-nearest}
\cindex{rounding!to nearest|hyperpagebf}%
Given \RSet{} the set of real numbers and \FSet{} the set of floating-point
numbers (\code{double}), $\roundNearest{x}$ is, for $x\in\RSet$, the
$y\in\FSet$ that minimizes $|x-y|$:
\begin{equation*}
  |x-\roundNearest{x}| = \min\{|x-y|\mid y\in\FSet\},
\end{equation*}
the one whose significand is even when two are equally near; a real of magnitude
at least $2^{1024}-2^{970}$ is rounded to the infinity of its sign, as IEEE 754
does.
\end{Def}

For a double $d$, the three roundings return $d$ itself: a double is never
moved. For a real $x$ that is no double, $\roundDn{x}$ and $\roundUp{x}$
are two consecutive doubles, and $\roundNearest{x}$ is one of them. Beyond
the largest double, the definitions give $\roundUp{x}=+\infty$ and
$\roundDn{x}=\mathrm{MAX}$ for $x>\mathrm{MAX}$.

\begin{Def}[Neighbours of a double]\label{def:neighbours}
\cindex{double!neighbours}%
For a double $d$, $\nextdouble{d}=\min\{y\in\FSet\mid y>d\}$ is the double
that follows $d$, and $\prevdouble{d}=\max\{y\in\FSet\mid y<d\}$ the double
that precedes it. For a real $x$ that is no double,
$\roundUp{x}=\nextdouble{(\roundDn{x})}$. The function \code{next\_float()}
of GAOL computes the successor of a double (Chapter~\ref{chap:errors}).
\end{Def}

\begin{Ex}\label{ex:one-tenth}
Programming with floating-point numbers is one of the few activities where
the compiler has to be considered with distrust. Let us consider the
following piece of code:
\begin{example}
#include <gaol/gaol>  @rem^// We do not import the gaol namespace~

using gaol::interval;

int main(void) {
  interval one_tenth(0.1); @rem^// BEWARE: wrong!~

  @emph^[some code using one_tenth]~
  gaol::cleanup();
}
\end{example}
Though a rational perfectly representable in decimal, \code{0.1} is not
representable in binary, at least not with a finite number of bits, and
thus requires rounding. The purpose of the programmer was to define an
interval containing this value. However, \code{0.1} is rounded \emph{at
compile time}, most certainly to the nearest representable floating-point
number $\roundNearest{0.1}$. As a consequence, \code{one\_tenth} is a
degenerate interval containing only $\roundNearest{0.1}$, and the
containment property (\exref{sec:containment}) is violated. The right way
to write such a constant is given in Example~\ref{ex:one-tenth-right}.
\end{Ex}

\section{Intervals}\label{sec:intervals}
%%-------------------
\cindex{interval!floating-point}\cindex{I@\ISet}%

\begin{Def}[Floating-point interval]\label{def:interval}
A \dfn{floating-point interval} $I=\itv{a}{b}$, with $a$ and $b$ in \FSet,
$a\leq b$, $a<+\infty$ and $b>-\infty$, is the set of the real numbers
between its bounds:
\begin{equation*}
  \itv{a}{b}=\{x\in\RSet\mid a\leq x\leq b\}.
\end{equation*}
$a=\leftBound{I}$ is its \dfn{lower bound}, or left bound, and
$b=\rightBound{I}$ its \dfn{upper bound}, or right bound. The empty set
$\emptyset$ is an interval as well. We write \ISet{} for the set of the
floating-point intervals.
\end{Def}

An interval is a set of real numbers: an infinite bound marks that the
interval is unbounded on that side, and is no member of it. The intervals
with an infinite bound, as \itv{1}{+\infty}, are the \dfn{unbounded}
intervals, intervals as any other; \itv{-\infty}{+\infty} is the whole real
line, which IEEE 1788-2015 calls \ieee{entire} and GAOL
\code{interval::universe()}. The point intervals $\itv{+\infty}{+\infty}$
and $\itv{-\infty}{-\infty}$ do not exist: the standard has no such
interval, and the constructors of GAOL v5 give the empty set for them
(\pxref{par:inf_inf_note}).\newinvfive{} An interval $\itv{a}{a}$ reduced to
one double is a \dfn{degenerate} interval, or point interval, which the
literals of the standard write $[a]$. The empty set is the result of an
operation whose argument has no point in the domain of the function
(\exref{sec:ieee1788-overview}); GAOL holds it as an interval whose two
bounds are NaN, and the method \code{is\_empty()} tells it.
\cindex{interval!unbounded}\cindex{interval!degenerate}\cindex{empty set}%

\begin{Def}[Hull]\label{def:hull}
\cindex{hull}%
The \dfn{hull} of a set $S$ of real numbers, written $\hull{S}$, is the
smallest floating-point interval containing $S$: $\hull{\emptyset}=
\emptyset$, and for a non-empty $S$
\begin{equation*}
  \hull{S}=\itv{\roundDn{\inf S}}{\roundUp{\sup S}},
\end{equation*}
where $\inf S=-\infty$ when $S$ has no lower bound and $\sup S=+\infty$ when
it has no upper bound.
\end{Def}

\begin{Def}[Canonical interval]\label{def:canonical}
\cindex{interval!canonical|see{canonical interval}}%
\cindex{canonical interval|hyperpagebf}%
A non-empty interval $I=\itv{a}{b}$ is \dfn{canonical} if and only if
$a\geq\prevdouble{b}$: it contains at most two consecutive floating-point
numbers.
\end{Def}

The hull of a real number $x$ that is no double,
$\hull{\{x\}}=\itv{\roundDn{x}}{\roundUp{x}}$, is canonical: it is the
tightest interval of doubles that encloses $x$, and no element of \ISet{}
does better.

\begin{Ex}\label{ex:one-tenth-right}
The right way to deal with rational constants that might not be perfectly
representable as floating-point numbers is to write them as strings, and to
read them with \code{gaol::textToInterval()} (\pxref{sec:constructors}), so
that they are correctly rounded downward and upward at \emph{run time}:
\begin{example}
#include <gaol/gaol>

using gaol::interval;

int main(void)
{
  @rem^// OK: this is the right way~
  interval one_tenth = gaol::textToInterval("0.1");
  @rem^// Another possible way~
  interval one_tenth2 = gaol::textToInterval("1/10");

  @emph^[some code using one_tenth]~
  gaol::cleanup();
  return 0;
}
\end{example}
Now, \code{one\_tenth} is the smallest floating-point interval enclosing
$0.1$: a canonical interval. Written in hexadecimal
(\pxref{sec:output-format}), which shows the bounds exactly, the interval of
Example~\ref{ex:one-tenth} and this one are:
\begin{onscreen}
interval(0.1) = [0x1.999999999999ap-4]
one_tenth     = [0x1.9999999999999p-4, 0x1.999999999999ap-4]
\end{onscreen}
\end{Ex}

\paragraph{Width, midpoint and radius.}
Three numbers describe an interval $\itv{a}{b}$: its \dfn{width} $b-a$
(\ieee{wid} of IEEE 1788-2015), its \dfn{midpoint} $(a+b)/2$ (\ieee{mid})
and its \dfn{radius} $(b-a)/2$ (\ieee{rad}), so that $\itv{a}{b}$ is the
set of the reals at a distance at most the radius from the midpoint.
Computed on doubles, none of the three is exact in general: the width GAOL
computes is rounded upward, and the midpoint is a double that the
computation has to keep inside the interval, for an unbounded interval
too~\cite{Goualard:TOMS2014}. Chapter~\ref{chap:functions} gives the
methods \code{width()}, \code{midpoint()} and \code{rad()}, and their
values on the empty set and on the unbounded intervals.
\cindex{interval!width}\cindex{interval!midpoint and radius}%

\section{Interval arithmetic and the containment property}
\label{sec:containment}
%%--------------------------------------------------------
\cindex{containment property}\cindex{outward rounding}%
\cindex{rounding!outward|see{outward rounding}}%

An operation on intervals stands for the operation on every pair of points
of its arguments. Its result has to contain all the values the real
operation can take and, to be of use, should contain as little else as
possible: the smallest floating-point interval that does is the hull.

\begin{Def}[Containment property]\label{def:containment}
The \dfn{containment property} of floating-point interval arithmetic
requires that, for any operation ``$\circ$'' and any intervals $I$ and $J$,
$I\circ J$ contains $\{i\circ j\mid i\in I,\ j\in J\}$, the pairs on which
$\circ$ is not defined being left out, and, for any function $f$ of one
variable and any interval $I$, that $f(I)$ contains $\{f(i)\mid i\in I\}$,
the points of $I$ outside the domain of $f$ being left out. The
\dfn{containment principle} asks for the smallest interval that does so, the
hull:
\begin{equation*}
  I\circ J=\hull{\{i\circ j\mid i\in I,\ j\in J\}},\qquad
  f(I)=\hull{\{f(i)\mid i\in I\}}.
\end{equation*}
\end{Def}

\begin{Ex}\label{ex:sqrt-functional}
For the square root, the principle gives
$\sqrt{I}=\hull{\{\sqrt{i}\mid i\in I\}}$. From monotonicity
considerations, the square root of \itv{1}{2} is then
$\sqrt{\itv{1}{2}}=\itv{1}{\sqrt{2}}$ on the real numbers; as $\sqrt2$ is
no double, the interval computed is $\itv{1}{\roundUp{\sqrt{2}}}$, whose
upper bound is the double above $\sqrt2$.
\end{Ex}

Fulfilling the containment property with doubles requires \dfn{outward
rounding}: the lower bound of a result is rounded downward and its upper
bound upward. For example, to add the intervals \itv{a}{b} and \itv{c}{d},
we compute $\itv{\roundDn{a+c}}{\roundUp{b+d}}$, where $\roundDn{r}$ and
$\roundUp{r}$ return the greatest double at most $r$ and the smallest double
at least $r$ (Definition~\ref{def:round-down-up}). These two operations are
performed by switching the rounding direction of the floating-point unit
towards, respectively, $-\infty$ and $+\infty$. With
$P=\{ac,\,ad,\,bc,\,bd\}$ and $Q=\{a/c,\,a/d,\,b/c,\,b/d\}$, the four basic
operations on bounded intervals become
\begin{align*}
  \itv{a}{b}+\itv{c}{d} &= \itv{\roundDn{a+c}}{\roundUp{b+d}},\\
  \itv{a}{b}-\itv{c}{d} &= \itv{\roundDn{a-d}}{\roundUp{b-c}},\\
  \itv{a}{b}\times\itv{c}{d} &= \itv{\roundDn{\min P}}{\roundUp{\max P}},\\
  \itv{a}{b}\div\itv{c}{d} &= \itv{\roundDn{\min Q}}{\roundUp{\max Q}}
    \quad\text{if } 0\notin\itv{c}{d}.
\end{align*}
A divisor that contains zero gives the hull of the quotients by its
non-zero points, which IEEE 1788-2015 defines case by case: an unbounded
interval when the dividend is not \itv{0}{0}, as
$\itv{1}{2}\div\itv{-1}{1}=\itv{-\infty}{+\infty}$, the interval \itv{0}{0} when
it is, and the empty set when the divisor is $\itv{0}{0}$, as
$\itv{1}{2}\div\itv{0}{0}=\emptyset$. The infinite bounds and the zeros of
each operation follow the standard as well; Chapter~\ref{chap:arithmetic}
states the rules.

\begin{Prop}\label{prop:round-neg}
For every real number $r$, $\roundDn{-r}=-\roundUp{r}$.
\end{Prop}
\demo The set \FSet{} is symmetric: $y\in\FSet$ if and only if
$-y\in\FSet$. Hence $\{y\in\FSet\mid y\leq-r\}=\{-z\mid z\in\FSet,\ z\geq
r\}$, and the greatest element of the first set is the opposite of the
smallest element of the second:
$\max\{y\in\FSet\mid y\leq-r\}=-\min\{z\in\FSet\mid z\geq r\}$, that is,
$\roundDn{-r}=-\roundUp{r}$.\boite

\medskip
On most platforms, switching the rounding direction is costly.
Proposition~\ref{prop:round-neg} cuts down the number of switches: nearly
all downward rounding operations can be replaced by upward rounding ones by
negating appropriately twice the operations performed. The lower bound of a
sum, $\roundDn{a+c}$, is $-\roundUp{(-a)+(-c)}$: the sum of the negated
lower bounds, rounded upward, then negated. GAOL stores the lower bound of
an interval negated, as the pair $\langle -a, b\rangle$, and computes both
bounds rounding upward: the rounding direction can then be set towards
$+\infty$ once and for all. The sum of $\langle -a, b\rangle$ and
$\langle -c, d\rangle$ is the pair $\langle\roundUp{-a-c},
\roundUp{b+d}\rangle$, two upward additions, which the kernels of GAOL
compute on the two halves of one register (\pxref{sec:kernels}). How GAOL
sets the rounding direction, and what the upward rounding does to the rest
of the program, is the subject of Chapter~\ref{chap:rounding}
(\pxref{sec:rounding-direction}).
\cindex{rounding!direction}\cindex{lower bound!stored negated}%

\paragraph{Exact operations.} Not every operation needs a rounding. The
negation maps \itv{a}{b} to \itv{-b}{-a}, whose bounds are doubles since
\FSet{} is symmetric; the minimum and the maximum return a bound of an argument,
and the absolute value a bound of its argument, its opposite or $0$; and the
integer functions, floor,
ceiling, truncation, the roundings to the nearest integer and the sign,
return integers that are doubles, the integer part of a double being a
double. These operations are exact: their result is the hull of the exact
result with no rounding at all, whatever the rounding direction in effect.
The four basic operations, the square root and the elementary functions
round: their bounds need a rounding in a chosen direction, and, for the
elementary functions, depend on the mathematical library that computes them
(\exref{sec:math-libraries}).
\cindex{exact operations}%

\section{The dependency problem}\label{sec:dependency}
%%--------------------------------
\cindex{dependency problem}%

The containment principle makes each operation as tight as doubles allow,
but a sequence of operations is not: an expression evaluated on intervals
gives an interval that encloses the range of the expression, and is in
general wider than it. The reason is that each occurrence of a variable
ranges over its interval on its own, as if the other occurrences did not
exist. This is the \dfn{dependency problem} of interval
arithmetic~\cite{Moore:Book1966,Neumaier:Book1990}.

\begin{Ex}\label{ex:dependency}
Let $x=\itv{0}{1}$. The expression $x-x$ is zero for every value of $x$,
but the subtraction of intervals takes its two operands as independent:
$\itv{0}{1}-\itv{0}{1}=\itv{-1}{1}$. In the same way, with $x=\itv{-1}{1}$,
the product $x\times x=\itv{-1}{1}$ combines the value $-1$ of the first
factor with the value $1$ of the second, which one number cannot take at
once, whereas the square, which knows that its two factors are equal, gives
$x^2=\itv{0}{1}$, the range; GAOL provides it as \code{sqr()}
(Chapter~\ref{chap:arithmetic}). The distributivity is weakened in the same
way: for any intervals $x$, $y$ and $z$, the operations being computed
exactly on the reals,
$z\times(x+y)\subseteq z\times x+z\times y$, with equality not in general,
as $\itv{1}{2}\times(\itv{-1}{1}+\itv{3}{4})=\itv{2}{10}$ where
$\itv{1}{2}\times\itv{-1}{1}+\itv{1}{2}\times\itv{3}{4}=\itv{1}{10}$. Computed
with doubles and rounded outward, even the inclusion may fail by one double.
\end{Ex}

The effect grows with the expression. The first program of
Chapter~\ref{chap:overview} (\pxref{sec:first-program}) evaluates the
Goldstein-Price function on the two intervals \itv{-2}{2} of its variables,
and obtains an enclosure of its range that is 231 times as wide as the range:
each occurrence of $x$ and of $y$ in the formula ranges over \itv{-2}{2} on
its own, as if the others did not exist.
\cindex{enclosure}%

\medskip
Cutting $x$ and $y$ into smaller intervals, and taking the hull of the
values of $f$ over the pieces, narrows the enclosure, at the price of more
evaluations: the program \file{examples/03\_dependency\_problem.cpp} of the
sources does it. Writing the expression with fewer occurrences of each
variable, the square in place of a product of a variable by itself for
instance, is the other remedy: the width of an enclosure is a property of
the expression evaluated, not of the function alone.

\section{IEEE 1788-2015}\label{sec:ieee1788-overview}
%%-----------------------
\cindex{IEEE 1788-2015}%

IEEE 1788-2015~\cite{IEEE1788:2015} is the standard for interval
arithmetic. It fixes what the libraries had each decided their own way:
what an interval is, what the operations return at the bounds of their
domains, how the empty set, the infinite bounds and the comparisons behave,
how an interval is written and read, and how tight each operation has to
be. GAOL v5 follows the standard in the operations GAOL provides: the empty
set, the infinite bounds, the domains of the functions, the comparisons and
the interval literals. The few cases where it differs, mostly the
\code{pow} of the namespace \code{gaol}, are stated where the operations
are described (\pxref{sec:pow-namespaces}, Chapter~\ref{chap:ieee1788}),
and the tightness of each operation is documented, as the standard requires
(Chapter~\ref{chap:accuracy}).

\paragraph{The set-based flavor and the bare intervals.} The standard is
organized in \emph{flavors}, and defines one, the \dfn{set-based flavor}:
an interval is a closed connected subset of \RSet, that is, the empty set,
a bounded interval \itv{a}{b} or an unbounded one, and the exact result of
an operation is the set of its values on the points of its arguments that
lie in the domain of the function, which the interval returned encloses
(\exref{sec:accuracy-modes}). The intervals of GAOL are those of the
set-based flavor (\exref{sec:intervals}). The standard also attaches to
its intervals a \emph{decoration}, which records whether the function was
defined and continuous on the whole of its argument along the computation
(Clause 11). GAOL has no decorations: only bare intervals are provided, and
the other parts of the standard it does not provide are listed in
Chapter~\ref{chap:ieee1788} (\pxref{sec:ieee1788-missing}).
\cindex{decorations}\cindex{bare interval}\cindex{set-based flavor}%

\paragraph{The operations.} The standard lists its operations in tables.
Table~9.1 holds the forward elementary functions every implementation
provides, with the domain of each: the arithmetic operations, the powers,
the exponentials and logarithms, the trigonometric and hyperbolic functions
and their inverses, the integer functions, \ieee{abs}, \ieee{min} and
\ieee{max}. Table~10.1 of the standard holds the reverse functions
(\exref{sec:relational-idea}), Table~10.2 the numeric functions of an
interval, Tables~10.3 and 10.4 the comparisons, and Table~10.5 the
recommended functions, which an implementation should provide (10.6). GAOL v5
adds the operations of these tables that GAOL~4 lacked, which the standard
requires or recommends, and the introduction lists them.\newinvfive{} The
relations of Tables~10.3 and 10.4, \ieee{equal}, \ieee{subset},
\ieee{interior}, \ieee{disjoint}, \ieee{precedes}, \ieee{strictPrecedes},
\ieee{less} and \ieee{strictLess}, follow the definitions of the standard
for the empty set and for the infinite bounds too
(Chapter~\ref{chap:relations}).\newinvfive{} The table of
Chapter~\ref{chap:ieee1788} says which operation of the standard each
function of GAOL is.
\cindex{IEEE 1788-2015!relations}%

\paragraph{The conventions.} A few conventions of the standard run through
the whole manual.
\begin{itemize}
\item[$\bullet$] \textbf{The domains.} A function is computed on the part
  of its argument within its domain, which is that of IEEE 1788-2015
  (Tables~9.1 and 10.5): the part outside is left out, and an argument with
  no point of the domain gives the empty set, as $\code{log}(\itv{-4}{0})$
  does.\newinvfive{} Every function gives the empty set for an empty
  argument. GAOL, which has no decorations, says nothing of the part left
  out: $\code{sqrt}(\itv{-1}{4})$ is \itv{0}{2}.
  \cindex{domain of a function}%
\item[$\bullet$] \textbf{The empty set.} The empty set is an interval, the
  result of an operation on an argument outside the domain of the function,
  and an interval whose bounds are in the wrong order, as
  \code{interval(1, 0)}, is empty too. Its numeric functions are those of
  Table~10.2: \ieee{inf} and \ieee{sup} of the empty set are $+\infty$ and
  $-\infty$, where \code{left()} and \code{right()} of GAOL are NaN, and
  \ieee{inf} returns $-0$ for a left bound zero (12.12.8).
\item[$\bullet$] \textbf{The infinite bounds.} An interval is a set of
  real numbers: an infinite bound is no member of it, so that
  \ieee{isMember}$(m, x)$ is \code{false} for an infinite $m$ (10.6.3), and
  the point intervals at $\pm\infty$ do not exist. The relations follow
  these definitions for the infinite bounds too.
\item[$\bullet$] \textbf{The literals.} An interval is written as a
  \dfn{literal} of the standard (12.11): \code{[l, r]}, \code{[a]} for a
  degenerate interval, \code{[empty]} for the empty set and \code{[entire]}
  for the whole line, the case of the letters being ignored; a literal is
  read back as an interval containing the one written (13.2, 13.3). GAOL writes
  and reads them (Chapter~\ref{chap:io}, \pxref{sec:input-format});
  \code{textToInterval} of the namespace \code{gaol\_ieee1788} reads the
  names of the functions of the standard and returns the empty set for a
  string that is no interval (12.1.3), where that of \code{gaol} reads the
  names of GAOL and throws (\pxref{sec:ieee1788-text}).\newinvfive
  \cindex{interval literal}%
\item[$\bullet$] \textbf{The powers.} The \ieee{pow} of Table~9.1 is
  defined for $x>0$, and for $x=0$ when $y>0$; the power with an integer
  exponent is \ieee{pown}$(x, n)$, whose $n$ is an integer rather than an
  interval. The \code{pow} of the namespace \code{gaol} takes the integer
  power for an integer exponent, a negative base included, which is the
  main point where GAOL and the standard differ (\pxref{sec:pow-namespaces},
  Section~\ref{sec:powers}).
\end{itemize}

\medskip
\noindent GAOL names its operations its own way, and the names of the
standard differ: the reverse function \ieee{coshRev}$(c, x)$ is
\code{acosh\_rel(c, x)}, and \ieee{roundTiesToEven}$(x)$ is
\code{round\_ties\_to\_even(x)}. The type \code{interval} and the functions
of GAOL live in the namespace \code{gaol}, and the namespace
\code{gaol\_ieee1788}, which \file{gaol/gaol.h} brings along, gives each
operation of the standard that GAOL provides the name and the argument
order the standard gives it, so that a reader of the standard finds it
without looking for its translation (\pxref{sec:namespaces},
Chapter~\ref{chap:ieee1788}).\newinvfive
\cindex{namespace!gaol\_ieee1788}\cindex{IEEE 1788-2015!names}%

\section{Relational (reverse) functions}\label{sec:relational-idea}
%%--------------------------------------
\cindex{operator!relational}\cindex{relational operator}%
\cindex{relational function}\cindex{reverse function}%

A function on intervals, as the containment principle defines it, maps an
interval of arguments to the hull of its values: this is its
\dfn{functional} version. A second reading of the same function is of use
when the function appears in a relation between variables.

\begin{Ex}\label{ex:sqrt-relational}
The functional square root of \itv{1}{2} is
$\sqrt{\itv{1}{2}}=\itv{1}{\sqrt{2}}$ (Example~\ref{ex:sqrt-functional}).
Another interpretation of the square root is
\begin{equation*}
  \sqrt{_r\,I}=\hull{\{j\in\RSet\mid\exists i\in I\colon j^2=i\}},
\end{equation*}
the \dfn{relational square root}, which permits obtaining both negative and
positive values: $\sqrt{_r\,\itv{1}{2}}=\itv{-\sqrt{2}}{\sqrt{2}}$. This
operator arises when we consider the relation $x^2=y$, which can
alternatively be written $x=\sqrt{_r\,y}$. Here, the \dfn{functional square
root} is not suitable, since it would confine $x$ to $\itv{0}{+\infty}$ and
lose the negative solutions of the relation.
\end{Ex}

Some applications, mainly in the area of constraint programming, require
such operators: a contractor narrows the domains of the variables of a
relation, and needs, for the relation $y=f(x)$, the part of the domain of
$x$ whose image meets the domain of $y$. As a consequence, GAOL offers both
functional and relational versions of the main arithmetic operators, which
IEEE 1788-2015 has since required, as its \dfn{reverse functions}
(Table~10.1).

\medskip
A relational function of one argument, as the relational square root, is
enough when the preimage of an interval has finitely many pieces. For a
periodic function the preimage has infinitely many, and the domain of the
variable has to be taken into account.

\begin{Ex}\label{ex:acos-relational}
Consider the relation $y=\cos x$, where $x$ and $y$ are interval variables.
One would like to express this relation in the equivalent way
$x=\operatorname{acos} y$. However, the function $\operatorname{acos}$
cannot be used, because its result is always included in the interval
\itv{0}{\pi}. What is needed is a relational version of
$\operatorname{acos}$; but since, for any value $y$, there are infinitely
many values $x$ verifying $y=\cos x$, the domain of $x$ has to be taken
into account. GAOL therefore defines a binary operator \code{acos\_rel} as
\begin{equation*}
  \code{acos\_rel}(Y, X)=\hull{\{x\in X\mid\exists y\in Y\colon y=\cos x\}},
\end{equation*}
to be contrasted with the functional arccosine,
\begin{equation*}
  \code{acos}(Y)=\hull{\{x\in\RSet\mid\exists y\in Y\colon
    x=\operatorname{acos} y\}}.
\end{equation*}
Figure~\ref{fig:relational-cos}, in Chapter~\ref{chap:arithmetic}, presents
the different results obtained when computing $\code{acos}(J)$ or
$\code{acos\_rel}(J, I)$, where $J$ and $I$ play the parts of $Y$ and $X$.
An interval $X$ that meets no piece of the
preimage gives the empty set.\newinvfive
\end{Ex}

\begin{Def}[Relational function]\label{def:relational}
Let $f$ be a function of one real variable, $Y$ an interval of values and
$X$ an interval of arguments. The \dfn{relational function} of $f$ with
respect to $X$, its \dfn{reverse function} in the words of IEEE 1788-2015,
is
\begin{equation*}
  f_{\mathrm{rel}}(Y, X)=\hull{\{x\in X\mid\exists y\in Y\colon y=f(x)\}},
\end{equation*}
the hull of the part of $X$ whose image by $f$ meets $Y$. The relational
function of one argument is the case $X=\RSet$.
\end{Def}

GAOL names a relational function after the inverse of the function it
reverses, as \code{acos\_rel} for the cosine, \code{asin\_rel} for the sine
and \code{sqrt\_rel} for the square; the standard names it after the
function itself, \ieee{cosRev}, \ieee{sinRev}, \ieee{sqrRev}, and the
namespace \code{gaol\_ieee1788} gives these names
(Chapter~\ref{chap:ieee1788}). The relational version of a binary
operation, as the relational division \code{div\_rel}, which reverses the
product, takes one more argument. The reference of the relational
arithmetic, with the figure of the relational cosine, is in
Chapter~\ref{chap:arithmetic} (\pxref{sec:relational-arithmetic}).

\section{The elementary functions and the mathematical libraries}
\label{sec:math-libraries}
%%--------------------------------------------------------------
\cindex{elementary functions}\cindex{mathematical library}%
\cindex{CORE-MATH}\cindex{libm}\cindex{mathlib}\cindex{CRlibm}%

The bounds of the four basic operations and of the square root are the
business of the processor: IEEE 754 requires them correctly rounded in the
direction in effect, and the roundings of Definition~\ref{def:round-down-up}
are what the floating-point unit computes. The elementary functions,
$\exp$, $\log$, $\sin$, $\cos$, $\tan$, their inverses and the hyperbolic
functions, are computed by software, by a mathematical library, and the
tightness of their bounds depends on it. GAOL v5 bounds its
elementary functions with CORE-MATH~\cite{CoreMath:2022}, compiled into it
from \file{3rd/math-core}, and takes no elementary function from any other
library: the C library gives it only \code{sqrt}, \code{floor}, \code{fma} and
the functions of \file{fenv.h}, and the integer powers, the roots and
\code{compoundm1} are computed by GAOL itself. Before that, GAOL was built with
one of three: mathlib by default with its \cmd{configure}, CRlibm, which its
meson build took by default, or the mathematical library of the system, whose options the three builds of GAOL
v5 no longer have (Chapter~\ref{chap:changes}).\newinvfive{} The
four come from four moments of the same history and were written with
different goals, which is why they give bounds of different tightness. This
section says where each of them comes from and what separates them, and why
GAOL v5 ended up with the last.

\subsection{The table maker's dilemma}
\cindex{table maker's dilemma}\cindex{correct rounding}%
\cindex{faithful rounding}%

A mathematical library evaluates $\exp$, $\log$, $\sin$\ldots{} in a
format, binary64 here, in which the exact result is not representable: it
has to round a value it only knows approximately. For some arguments the
exact result falls extraordinarily close to the middle of two consecutive
doubles, or to a double for the directed roundings, and deciding which way to round then asks for far more than 53
bits of intermediate precision. No bound on that distance follows from the
definition of the functions: this is the \dfn{table maker's dilemma}, named
by W. Kahan. Two ambitions follow from it, and they are the axis along
which the four libraries differ.

\begin{Def}[Faithful rounding]\label{def:faithful}
A library is \dfn{faithfully rounded}, or ``nearly correct'', when its
error is below one unit in the last place: it returns one of the two doubles
next to the exact value. The libraries of this kind are within about half a
unit in practice. This is enough for most numerical computation.
\end{Def}

\begin{Def}[Correct rounding]\label{def:correct-rounding}
A library is \dfn{correctly rounded} when it returns exactly the double the
rounding direction in effect gives of the exact result:
$\roundNearest{f(x)}$, $\roundDn{f(x)}$ or $\roundUp{f(x)}$. It makes a
program reproducible from one machine and one library to another, and, for
interval arithmetic, it gives the tightest bound a format can hold: a value
correctly rounded upward \emph{is} the upper bound, with nothing to add.
\end{Def}

\noindent IEEE 754-2008 lists the correctly rounded elementary functions
among its recommended operations, which is what the last two libraries
below set out to provide.

\subsection{libm, the mathematical library of the system}

libm is not a library but an interface: the mathematical part of the C
standard library, with one implementation per system. The line starts with
the libm of Unix V7 in the 1970s, then the libm of 4.3BSD in the middle of
the 1980s, written around W. Kahan at Berkeley (Coonen, K. C. Ng, Thomas,
Tang), where the culture of testing the accuracy of elementary functions
was formed. In 1993 Sun published fdlibm, the Freely Distributable libm of
K. C. Ng: portable C, no assembly, aiming at less than one unit in the last
place. fdlibm is the common ancestor of nearly everything that followed:
\file{msun} of FreeBSD, the libm of musl, openlibm of Julia, and part of
that of the glibc.

\medskip
Its design puts speed first and accuracy second. There is one evaluation
phase, an argument reduction and a polynomial or rational approximation
chosen by hand for each function; the error bound announced is empirical,
established by tests rather than proved; the mathematical properties
(monotonicity, symmetry, $\sin^2x+\cos^2x\leq1$) are not guaranteed.
Rounding to nearest is assumed throughout: nothing is promised in the
directed rounding directions, which is precisely what interval arithmetic
needs. Beside these portable implementations live the fast and
non-portable ones: Intel SVML, AMD libm, the \file{optimized-routines} of
ARM, and the vector variants.

\medskip
This is why GAOL could not keep the libm of the system as an option: the
bounds would not be certified, holding only where that libm stays within
about one double of the exact value.

\subsection{mathlib, the IBM Accurate Portable Mathematical Library}

mathlib, also called libultim, was written at IBM Haifa at the end of the
1990s (Abraham Ziv, Moshe Olshansky, Ealan Henis, Anna Reitman). It is the
industrial form of the idea Ziv published in 1991 in \emph{ACM TOMS}: the
\dfn{onion peeling} strategy. The function is evaluated with increasing
precision, a fast phase accurate to about half a unit in the last place
plus a little, then a second one, then a software multiprecision
arithmetic, and the computation stops as soon as a rounding test proves
that the rounding can be decided. Since the hard cases are rare, the
average cost stays close to that of an ordinary libm, while the result is
correctly rounded in practice.

\medskip
Its two limits are the ones that decided its fate. Nothing bounds the
number of iterations: the strategy has no proof and no bound on the worst
case, which can cost thousands of cycles. And it rounds to nearest only, so
a bound drawn from it has to be moved one double outward, which costs one
unit in the last place of tightness: this is what GAOL~4 did, built with
mathlib, for every elementary function but the hyperbolic ones
(\exref{sec:accuracy-modes}).

\medskip
mathlib was contributed to the glibc around 2001 and provided its
\code{exp}, \code{log}, \code{sin}, \code{cos}, \code{atan} and \code{pow}
for some fifteen years. In 2018 the glibc 2.28 removed these
multiprecision slow paths: the pathological execution times were judged
worse than the accuracy was worth. The glibc thereby gave up correct
rounding and went back to about half to one unit in the last place with a
bounded worst case. The version GAOL built was mathlib 2.1.1, under the GNU
LGPL, as Fr\'ed\'eric Goualard distributed it; GAOL v5 carries it no
longer.\newinvfive

\subsection{CRlibm, the proofs}

CRlibm is a project of the LIP at the ENS de Lyon and then of the INRIA
Ar\'enaire team (later AriC), from 2000 to about 2011: Jean-Michel Muller,
Florent de Dinechin, Catherine Daramy-Loirat, David Defour, Christoph
Lauter, with Guillaume Melquiond for the proof tools. Its last release,
1.0beta4, dates from 2010; it is no longer maintained. CRlibm keeps Ziv's
two phases but changes what is guaranteed.
\begin{itemize}
\item[$\bullet$] \textbf{The second phase is bounded.} Lef\`evre and Muller
  had computed, by an exhaustive search made tractable by their algorithms,
  the hardest-to-round cases in double precision: about 120 bits are known
  to be enough. No unbounded multiprecision is needed any more, a
  double-double or triple-double arithmetic suffices, and the worst-case
  execution time is bounded, typically a few times the average.
\item[$\bullet$] \textbf{The error bounds are proved}, and in part checked
  by machine with Gappa. The documentation of CRlibm is as much a proof
  document as a manual, and this is its main scientific contribution.
\item[$\bullet$] \textbf{The four rounding directions are provided
  explicitly}: \code{cos\_rn}, \code{cos\_rd}, \code{cos\_ru},
  \code{cos\_rz}. This is what makes it natural for interval arithmetic,
  and why GAOL took its functions rounded downward and upward directly,
  without moving them outward: built with CRlibm, the \code{exp},
  \code{log}, \code{sin}, \code{cos}, \code{tan}, \code{asin}, \code{acos},
  \code{atan}, \code{sinh} and \code{cosh} of GAOL were the tightest bounds
  of the value at the bound. CORE-MATH gives the same, in the direction in
  effect, for every function and on every system
  (\exref{sec:accuracy-modes}).
\end{itemize}
Its price was the work of a specialist for each function, a narrow coverage
(a dozen functions, binary64 only: neither \code{tanh} nor \code{asinh},
\code{acosh}, \code{atanh}), and an LGPL licence that kept it out of the
libm of the systems. Its direct descendant is Metalibm (de Dinechin,
Kupriianova, Lauter): generating the code and its proof rather than writing
them.

\subsection{CORE-MATH, correct rounding made adoptable}

CORE-MATH~\cite{CoreMath:2022} started in 2021 at Inria (Caramba team,
Nancy), by Paul Zimmermann, also behind MPFR, with Alexei Sibidanov, Tom
Hubrecht, S\'el\`ene Corbineau and others. It draws the lessons of the
three libraries above, and its design is as much a strategy of adoption as
a technique.
\begin{itemize}
\item[$\bullet$] \textbf{The stated goal is correct rounding at least as
  fast as the current libms.} It answers the reason mathlib was removed
  from the glibc: if correct rounding costs nothing, the argument from
  latency falls.
\item[$\bullet$] \textbf{The licence is MIT}, deliberately permissive, so
  that the code can be taken up by the glibc, the libm of LLVM or musl. It
  answers the reason CRlibm was not adopted, and it works: routines of
  CORE-MATH have been integrated upstream. It is also what lets GAOL v5
  take the sources it needs as they are, in \file{3rd/math-core}, and compile
  them into its library.
\item[$\bullet$] \textbf{The coverage is wide}: the functions of C23, in
  binary32, binary64 and binary128, the recent functions included.
\item[$\bullet$] \textbf{The validation is graded}: in binary32 it is
  exhaustive, the $2^{32}$ arguments being enumerated and compared with
  MPFR; in binary64 and beyond, a proved fast phase and an accurate phase,
  with a search for the hard cases.
\item[$\bullet$] \textbf{The four rounding directions are honoured}, the
  one in effect included. This is what GAOL v5 rests on: correctly rounded
  in the upward rounding GAOL computes in, CORE-MATH gives the tightest
  bounds without switching the rounding direction: the interval \code{log} of
  GAOL v5 takes 20.9\,ns where that of GAOL~4.2.3, which calls mathlib, moves
  its value outward and switches the rounding direction twice, takes 46.8\,ns
  (Intel i7-1185G7, Clang 21.1, 17 October 2026, \pxref{sec:comparison}).
\end{itemize}
A proposal to add \code{cr\_*} functions to the C standard comes from the
same team, so that correct rounding can be asked for by the program rather
than imposed on everyone. Beside it, and of an orthogonal design, is RLIBM
(Rutgers, Jay Lim and Santosh Nagarakatte, about 2019--2022): instead of
approximating the real function and then rounding, it looks by linear
programming for a polynomial that directly yields the correct rounding for
every input. Elegant, but practicable mainly for the short formats
(binary32, bfloat16, posits).

\subsection{What separates them}

Table~\ref{tab:math-libraries} sets the four libraries side by side. The
line of the story is plain: libm established that a mathematical library
had to be portable and fast; mathlib showed that correct rounding was
reachable on average, and was rejected on its worst case; CRlibm showed
that it could be proved and bounded, and was not adopted for want of a
licence, of coverage and of maintenance; CORE-MATH takes all of it up and
treats industrial adoption, that is, performance, MIT licence, C23 coverage
and standardisation, as a design constraint of the same rank as accuracy.

\begin{table}[htbp]
  \centering\footnotesize\renewcommand{\arraystretch}{1.15}
  \setlength{\tabcolsep}{4pt}
  \caption{The four mathematical libraries GAOL has been built with} \label{tab:math-libraries}
  \begin{tabular}{L{1.6cm}L{2.8cm}L{2.5cm}L{2.9cm}L{2.9cm}}
    \hline
    & libm (fdlibm, glibc) & mathlib (IBM) & CRlibm & CORE-MATH\\
    \hline\hline
    Time & 1980s, 1993-- & about 1997--2001 & 2000--2011 & 2021--\\
    \hline
    Where from & Berkeley, Sun & IBM Haifa & ENS de Lyon, INRIA
      & INRIA, CERN\\
    \hline
    Goal & speed, below 1 ulp & correct rounding in practice
      & correct rounding proved & correct rounding and speed\\
    \hline
    Method & one phase & unbounded onion peeling
      & two phases, the second bounded by the hard cases
      & proved phases, exhaustive in binary32\\
    \hline
    Evidence & tests & tests & Gappa, published proofs
      & proofs and exhaustive checks\\
    \hline
    Worst case & bounded & unbounded & bounded & bounded\\
    \hline
    Rounding & to nearest & to nearest & the four, explicit
      & the four, the one in effect\\
    \hline
    Formats & all & binary64 & binary64 & binary32, 64, 128\\
    \hline
    Licence & permissive & LGPL & LGPL & MIT\\
    \hline
    State & alive & frozen & abandoned & active, being adopted\\
    \hline
  \end{tabular}

\end{table}

\subsection{What it means for the bounds of GAOL}

The hierarchy is read directly in the tightness of the bounds, which
Chapter~\ref{chap:accuracy} states function by function.
\begin{itemize}
\item[$\bullet$] \textbf{The mathematical library of the system} gave
  bounds that were not certified: the value rounded to nearest moved outward
  by $2\cdot2^{-52}$ relatively and by the smallest normal double, several
  doubles beyond the tightest one.
\item[$\bullet$] \textbf{mathlib} was correctly rounded to nearest, hence
  one double outward: the bounds were accurate within one double.
\item[$\bullet$] \textbf{CRlibm and CORE-MATH} are correctly rounded in the
  direction wanted, hence give the tightest bound of the value at the
  bound, with nothing to add.
\end{itemize}
\noindent This is also why GAOL v5 takes every one of its elementary
functions from CORE-MATH, above all the hyperbolic ones (neither mathlib nor
CRlibm has them all, and the libm of the system is not accurate enough
everywhere for its values moved outward to be bounds), and on every
architecture and with every compiler, \file{gaol/math/gaol\_u128.h} giving
the accurate phases the 128-bit integer where the compiler has none. The
bounds are the same on every architecture and with every compiler, and
there is no mathematical library to install, to find or to link along with
GAOL: CORE-MATH is compiled into it.\newinvfive

\paragraph{Further reading.} The onion peeling strategy is in A. Ziv,
\emph{Fast evaluation of elementary mathematical functions with correctly
rounded last bit}, ACM TOMS 17(3), 1991. The table maker's dilemma, the
hard cases and the two-phase evaluation are in J.-M. Muller,
\emph{Elementary Functions: Algorithms and Implementation}, and in the
\emph{Handbook of Floating-Point Arithmetic}; the floating-point arithmetic
itself is presented by Goldberg~\cite{Goldberg:01} and in the
standard~\cite{IEEE754:2008}. The documentation distributed with the sources
of CRlibm holds the proofs of its error bounds. CORE-MATH is presented
in~\cite{CoreMath:2022} and at \url{https://core-math.gitlabpages.inria.fr/};
the file \file{3rd/README.md} of the sources says what GAOL v5 takes from
it. The NEWS of the glibc 2.28 records the removal of the multiprecision
slow paths inherited from mathlib.

\section{The accuracy modes}\label{sec:accuracy-modes}
%%-----------------------------
\cindex{accuracy!modes}\cindex{tightest}\cindex{accurate}\cindex{valid}%
\cindex{tightness}%

The containment principle asks for the hull of the exact result, which the
basic operations reach but which the elementary functions, computed by a
library, may not. IEEE 1788-2015 therefore defines three accuracy modes
(12.10.1), for an operation $f$, its exact value $f^*$ and an input box
$x$, and requires an implementation to document the tightness of each of
its interval operations (12.10.3).

\begin{Def}[Tightest]\label{def:tightest}
$f$ is \dfn{tightest} when $f(x)$ is the hull, in doubles, of $f^*(x)$:
each bound is the nearest double on its side of the exact bound.
\end{Def}

\begin{Def}[Accurate]\label{def:accurate}
$f$ is \dfn{accurate} when $f(x)$ encloses $f^*(x)$ and is within
\begin{equation*}
  \mathrm{nextOut}\bigl(f_{\mathrm{tightest}}(\mathrm{nextOut}(x))\bigr),
\end{equation*}
$\mathrm{nextOut}$ moving each bound one double outward: roughly, $f(x)$
is no wider than the tightest result for an input one double wider, plus
one double on each side.
\end{Def}

\begin{Def}[Valid]\label{def:valid}
$f$ is \dfn{valid} when $f(x)$ encloses $f^*(x)$.
\end{Def}

The standard requires the basic operations (\ieee{neg}, \ieee{add}, \ieee{sub},
\ieee{mul}, \ieee{div}, \ieee{recip}, \ieee{sqr}, \ieee{sqrt}, \ieee{fma}),
the integer functions and \ieee{abs}, \ieee{min} and \ieee{max} to be
tightest; the other elementary functions, the reverse functions and the
recommended functions to be valid, and recommends them to be accurate
(12.10.2). In this manual, ``within $k$ doubles'' means that each bound is
at most $k$ doubles beyond the tightest one.

\medskip
GAOL bounds its elementary functions with CORE-MATH, which is correctly
rounded in the rounding direction in effect. Computing in the upward
rounding it keeps, GAOL takes the value at a bound as the upper bound, with
nothing to add, and the double below it as the lower one. The upper bound is
the tightest one, and so is the lower bound: the double below the upper bound
is the tightest lower bound when the exact value is no double, and where the
exact value is a double, as for $\log(1)=0$ and $\sin(0)=0$, the operations
return that value itself rather than the double below it
(Remark~\ref{rem:exact-value}). Where an algorithm
evaluates the function at the bounds of the interval, the bounds of the
interval are therefore the tightest ones. Before GAOL v5, built with mathlib,
the default of its \cmd{configure}, the elementary functions but the hyperbolic
ones were correctly rounded to nearest and moved one double outward, which is
at most one double beyond the tightest bound: accurate, and not tightest. The
hyperbolic functions came from the mathematical library of the system, whose
values are sometimes further than one double from the exact ones
(Remark~\ref{rem:accuracy-before}).\newinvfive

\begin{Rem}\label{rem:exact-value}
The double below the correctly rounded upper value is the tightest lower
bound only when the exact value is no double. When it is, as for
$\log(1)=0$, the value correctly rounded upward is the exact value itself,
and an operation that moved it one double down would be accurate, not
tightest. The operations of GAOL return the exact value where it happens,
and Chapter~\ref{chap:accuracy} says where.
\end{Rem}

\medskip
Chapter~\ref{chap:accuracy} documents the tightness of every operation of
\code{interval}, from the algorithm of each operation, as the standard
requires; the tests of \file{tests} check each statement on random and
special arguments (Chapter~\ref{chap:tests}), and the page \file{doc/accuracy.md}
gives, for each, the algorithm and the largest distance the tests
found.\newinvfive{} The statements hold for GAOL as the three builds make
it, whatever the rounding direction the calling code left, on every
architecture and with every compiler: for the SSE2, the NEON and the FPU
intervals, and for the AVX-512 path, whose bounds are those of the SSE2
intervals but for the sign of a zero bound of a product with a double. They hold under the conditions that chapter states: the
code including the headers of GAOL compiled with the flags of interval
arithmetic (\pxref{sec:compiling}), and the floating-point exceptions
masked as a program starts (\pxref{sec:fp-exceptions}). The intervals of
floats are not covered.

\section{Discussion}\label{sec:concepts-discussion}

This chapter has set the vocabulary of the manual. A floating-point
interval is a set of real numbers between two doubles, the empty set and
the unbounded intervals included; the containment property requires every
operation to return an interval that encloses the exact result, and
outward rounding is what enforces it on doubles. The property
$\roundDn{-r}=-\roundUp{r}$ lets GAOL compute both bounds in the upward
rounding, with the lower bound stored negated: a choice that spares the
switches of the rounding direction, at the price of leaving the direction
upward for the rest of the program, which Chapter~\ref{chap:rounding}
examines. The dependency problem is the limit of the method: an enclosure
is as tight as the expression allows, not as the function does.

\medskip
Two choices of GAOL v5 follow from the notions of this chapter. The
standard IEEE 1788-2015 fixes the special cases, the empty set, the
infinite bounds, the domains and the literals, and GAOL v5 follows it in
the operations it provides, under two sets of names; it provides bare
intervals only, without the decorations of the standard. The elementary
functions are bounded with CORE-MATH, whose correct rounding in the
direction in effect gives the tightest bounds without a switch of the
rounding direction, where the libm of the system gave uncertified bounds
and mathlib bounds one double too wide; the three accuracy modes of the
standard, tightest, accurate and valid, are the words in which
Chapter~\ref{chap:accuracy} states what each operation achieves.

\medskip
The next chapter leaves the concepts for the software: how GAOL v5 is
obtained, configured, built and installed on each platform
(Chapter~\ref{chap:installation}).
